% ============================================================
% Notes: Calculus of Variations
% Topics: variational principle, Euler--Lagrange equations,
%         constraints (Lagrange multipliers)
% Source style: Arfken & Weber, Mathematical Methods for Physicists
% ============================================================

\section{Variational calculus (calculus of variations)}

\subsection{Basic variational problem}
A standard variational problem is to find a function $y(x)$ that extremizes
(functional is stationary):
\begin{equation}
J[y] = \int_{x_1}^{x_2} F(x, y(x), y'(x))\,dx.
\end{equation}
We consider variations
\begin{equation}
y(x)\to y(x)+\epsilon\,\eta(x),
\end{equation}
where $\eta(x)$ is an arbitrary smooth test function satisfying boundary conditions.

\subsection{Fixed-endpoint conditions}
If endpoints are fixed:
\begin{equation}
y(x_1)=y_1,\qquad y(x_2)=y_2,
\end{equation}
then the variation must satisfy
\begin{equation}
\eta(x_1)=\eta(x_2)=0.
\end{equation}

\subsection{First variation and stationarity}
The condition for an extremum is
\begin{equation}
\delta J = 0
\quad \text{for all admissible }\eta(x).
\end{equation}
Compute:
\begin{align}
\delta J
&= \left.\frac{d}{d\epsilon}J[y+\epsilon\eta]\right|_{\epsilon=0}\\
&= \int_{x_1}^{x_2}
\left(
\frac{\partial F}{\partial y}\eta
+
\frac{\partial F}{\partial y'}\eta'
\right)dx.
\end{align}

Integrate by parts the $\eta'$ term:
\begin{align}
\int_{x_1}^{x_2}\frac{\partial F}{\partial y'}\eta'\,dx
&=
\left[\frac{\partial F}{\partial y'}\eta\right]_{x_1}^{x_2}
-\int_{x_1}^{x_2}\frac{d}{dx}\left(\frac{\partial F}{\partial y'}\right)\eta\,dx.
\end{align}
For fixed endpoints, the boundary term vanishes, so
\begin{equation}
\delta J
=
\int_{x_1}^{x_2}
\left(
\frac{\partial F}{\partial y}
-
\frac{d}{dx}\frac{\partial F}{\partial y'}
\right)\eta\,dx.
\end{equation}

Since $\eta(x)$ is arbitrary, the integrand must vanish.

% ------------------------------------------------------------

\section{Euler equation (Euler--Lagrange equation)}

\subsection{Euler--Lagrange equation (1 function)}
\begin{equation}
\boxed{
\frac{\partial F}{\partial y}
-
\frac{d}{dx}\left(\frac{\partial F}{\partial y'}\right)=0.
}
\end{equation}

\subsection{Multiple dependent variables}
If $y_i(x)$, $i=1,\dots,n$:
\begin{equation}
J[\{y_i\}] = \int_{x_1}^{x_2} F(x, y_1,\dots,y_n, y_1',\dots,y_n')\,dx,
\end{equation}
then for each $i$:
\begin{equation}
\boxed{
\frac{\partial F}{\partial y_i}
-
\frac{d}{dx}\left(\frac{\partial F}{\partial y_i'}\right)=0.
}
\end{equation}

\subsection{Special case: $F$ does not depend explicitly on $y$}
If $\partial F/\partial y = 0$, then
\begin{equation}
\frac{d}{dx}\left(\frac{\partial F}{\partial y'}\right)=0
\quad\Rightarrow\quad
\frac{\partial F}{\partial y'}=\text{constant}.
\end{equation}

\subsection{Special case: $F$ does not depend explicitly on $x$ (first integral)}
If $\partial F/\partial x = 0$, then
\begin{equation}
\boxed{
F - y'\frac{\partial F}{\partial y'} = \text{constant}.
}
\end{equation}
(This is the variational analog of energy conservation.)

% ------------------------------------------------------------

\section{Variational problems with constraints}

Constraints restrict the admissible functions $y(x)$.

\subsection{Integral constraint (isoperimetric problem)}
Suppose we extremize $J[y]$ subject to
\begin{equation}
I[y] = \int_{x_1}^{x_2} G(x,y,y')\,dx = C,
\end{equation}
where $C$ is fixed.

\subsection{Method of Lagrange multipliers (functional form)}
Introduce a constant multiplier $\lambda$ and define the augmented functional
\begin{equation}
\tilde J[y] = \int_{x_1}^{x_2}\Big(F(x,y,y') + \lambda\,G(x,y,y')\Big)\,dx.
\end{equation}
Then the constrained extremum is obtained by solving
\begin{equation}
\delta \tilde J = 0
\end{equation}
together with the constraint $I[y]=C$.

Thus the Euler--Lagrange equation becomes
\begin{equation}
\boxed{
\frac{\partial}{\partial y}\Big(F+\lambda G\Big)
-
\frac{d}{dx}\left(\frac{\partial}{\partial y'}\Big(F+\lambda G\Big)\right)=0,
}
\end{equation}
with $\lambda$ determined from the constraint.

\subsection{Pointwise constraint}
If there is a constraint that must hold at each $x$:
\begin{equation}
g(x,y,y')=0,
\end{equation}
then $\lambda$ may need to be a function $\lambda(x)$ and we extremize
\begin{equation}
\tilde J[y]=\int_{x_1}^{x_2}\Big(F(x,y,y')+\lambda(x)\,g(x,y,y')\Big)\,dx.
\end{equation}
Euler--Lagrange gives coupled equations for $y(x)$ and $\lambda(x)$:
\begin{equation}
\frac{\partial}{\partial y}(F+\lambda g)
-\frac{d}{dx}\left(\frac{\partial}{\partial y'}(F+\lambda g)\right)=0,
\end{equation}
and the constraint $g(x,y,y')=0$.

\subsection{Multiple constraints}
For $m$ integral constraints
\[
\int G_\alpha\,dx=C_\alpha,\qquad \alpha=1,\dots,m,
\]
use
\begin{equation}
\tilde J=\int\left(F+\sum_{\alpha=1}^m\lambda_\alpha G_\alpha\right)dx.
\end{equation}

% ------------------------------------------------------------

\section{Boundary conditions and transversality (quick)}

\subsection{Fixed endpoints}
If $y(x_1),y(x_2)$ are fixed, boundary terms vanish automatically.

\subsection{Free endpoint (natural boundary condition)}
If $y(x_2)$ is free, then $\eta(x_2)$ is not forced to be zero.
From the boundary term
\[
\left[\frac{\partial F}{\partial y'}\eta\right]_{x_1}^{x_2},
\]
we obtain the natural condition at $x_2$:
\begin{equation}
\boxed{
\left.\frac{\partial F}{\partial y'}\right|_{x=x_2}=0
\quad \text{(if }y(x_2)\text{ free)}.
}
\end{equation}
Similarly for $x_1$.

% ------------------------------------------------------------

\section{High-yield summary}

\begin{itemize}
\item Functional:
\[
J[y]=\int_{x_1}^{x_2}F(x,y,y')\,dx.
\]
\item Variation: $y\to y+\epsilon\eta$, stationarity $\delta J=0$.
\item Euler--Lagrange:
\[
\frac{\partial F}{\partial y}
-\frac{d}{dx}\left(\frac{\partial F}{\partial y'}\right)=0.
\]
\item If $\partial F/\partial x=0$:
\[
F-y'\frac{\partial F}{\partial y'}=\text{const}.
\]
\item Integral constraint $\int G\,dx=C$:
\[
F\to F+\lambda G.
\]
\item Pointwise constraint $g(x,y,y')=0$:
\[
F\to F+\lambda(x)g.
\]
\end{itemize}
